\documentclass[a4paper,USenglish,cleveref,autoref,thm-restate]{lipics-v2021}
%This is a template for producing LIPIcs articles. 
%See lipics-v2021-authors-guidelines.pdf for further information.
%for A4 paper format use option "a4paper", for US-letter use option "letterpaper"
%for british hyphenation rules use option "UKenglish", for american hyphenation rules use option "USenglish"
%for section-numbered lemmas etc., use "numberwithinsect"
%for enabling cleveref support, use "cleveref"
%for enabling autoref support, use "autoref"
%for anonymousing the authors (e.g. for double-blind review), add "anonymous"
%for enabling thm-restate support, use "thm-restate"
%for enabling a two-column layout for the author/affilation part (only applicable for > 6 authors), use "authorcolumns"
%for producing a PDF according the PDF/A standard, add "pdfa"

%\pdfoutput=1 %uncomment to ensure pdflatex processing (mandatatory e.g. to submit to arXiv)

\hideLIPIcs  %uncomment to remove references to LIPIcs series (logo, DOI, ...), e.g. when preparing a pre-final version to be uploaded to arXiv or another public repository
\nolinenumbers %uncomment to disable line numbering

%\graphicspath{{./graphics/}}%helpful if your graphic files are in another directory

\usepackage{tikz}
\usetikzlibrary{positioning}
\usetikzlibrary{calc}
\usepackage{censor}

%\usepackage[firstpageonly=true]{draftwatermark}

\bibliographystyle{plainurl}% the mandatory bibstyle

\title{A Formal Proof of Complexity Bounds on Diophantine Equations} 

\author{Jonas Bayer}{Department of Pure Mathematics and Mathematical Statistics, University of Cambridge, UK}{jcb234@cam.ac.uk}{https://orcid.org/0009-0006-2500-3146}{}
%mandatory, please use full name; only 1 author per \author macro; first two parameters are mandatory, other parameters can be empty. Please provide at least the name of the affiliation and the country. The full address is optional. Use additional curly braces to indicate the correct name splitting when the last name consists of multiple name parts.

\author{Marco David}{Department of Physics, University of California, Berkeley, US}{mdv@berkeley.edu}{https://orcid.org/0000-0001-5673-5389}{}

\authorrunning{J. Bayer and M. David} %mandatory. First: Use abbreviated first/middle names. Second (only in severe cases): Use first author plus 'et al.'

\Copyright{Jonas Bayer and Marco David} %mandatory, please use full first names. LIPIcs license is "CC-BY";  http://creativecommons.org/licenses/by/3.0/

\ccsdesc[100]{Theory of computation $\rightarrow$ Logic and verification} %mandatory: Please choose ACM 2012 classifications from https://dl.acm.org/ccs/ccs_flat.cfm 

\keywords{Diophantine Equations, Hilbert's Tenth Problem, Isabelle/HOL} %mandatory; please add comma-separated list of keywords

\category{} %optional, e.g. invited paper

\relatedversion{} %optional, e.g. full version hosted on arXiv, HAL, or other respository/website
%\relatedversiondetails[linktext={opt. text shown instead of the URL}, cite=DBLP:books/mk/GrayR93]{Classification (e.g. Full Version, Extended Version, Previous Version}{URL to related version} %linktext and cite are optional

%%% TODO UPDATE
%\supplement{The formal proof is available in Isabelle's Archive of Formal Proofs.}
%\supplementdetails[cite=up-afp]{Formal Proof Development}{https://isa-afp.org/TODO}
\supplement{The formal proof is available online.}
\supplementdetails[cite=up-afp]{Formal Proof Development}{https://gitlab.com/hilbert-10/universal-pairs}


%\funding{(Optional) general funding statement \dots}%optional, to capture a funding statement, which applies to all authors. Please enter author specific funding statements as fifth argument of the \author macro.

\acknowledgements{%
This project would not exist without Yuri Matiyasevich who proposed the formalization of Hilbert's Tenth Problem and encouraged our exploration of Universal Pairs. We are deeply grateful to Dierk Schleicher for his unwavering support throughout, both on a personal level and through his mathematical insight. We thank Malte Haßler, Thomas Serafini and Simon Dubischar for their mathematical contributions to the project.
%
We are indebted to the group of students who contributed at various stages of the formalization. In particular, we thank Timothé Ringeard and Xavier Pigé for leading the formal proof of the Bridge Theorem; Anton Danilkin for his work on \texttt{poly\_extract} and further developments in multivariate polynomial theory; Mathis Bouverot-Dupuis, Paul Wang, Quentin Vermande, and Theo André for their formalization of coding theory; Loïc Chevalier, Charlotte Dorneich, and Eva Brenner for the formalization of Matiyasevich’s polynomial and relation combining; and Chris Ye and Kevin Lee for their contributions to formal degree calculations.
%
The above highlights their primary contributions, though many have supported other aspects of this work whenever they were able.
We also thank all other members of the student work group for their involvement in the project. We gratefully acknowledge support from the D\'epartement de math\'emathiques et applications at \'Ecole Normale Sup\'erieure de Paris.
}



%Editor-only macros:: begin (do not touch as author)%%%%%%%%%%%%%%%%%%%%%%%%%%%%%%%%%%
\EventEditors{John Q. Open and Joan R. Access}
\EventNoEds{2}
\EventLongTitle{42nd Conference on Very Important Topics (CVIT 2016)}
\EventShortTitle{CVIT 2016}
\EventAcronym{CVIT}
\EventYear{2016}
\EventDate{December 24--27, 2016}
\EventLocation{Little Whinging, United Kingdom}
\EventLogo{}
\SeriesVolume{42}
\ArticleNo{23}
%%%%%%%%%%%%%%%%%%%%%%%%%%%%%%%%%%%%%%%%%%%%%%%%%%%%%%

%%%%%%
%%%%%% CUSTOM MACROS
%%%%%%
\usepackage[normalem]{ulem}
\usepackage{needspace}

\newcommand{\reminder}[1]{\textsl{{\color{blue} #1}}\protect\marginpar{$\lhd\lhd$}}
\newcommand{\Reminder}[1]{\textsl{#1}\protect\marginpar{$\lhd\lhd\lhd$}}

\newcommand{\N}{{\mathbb N}}
\newcommand{\Z}{{\mathbb Z}}
\newcommand{\B}{\mathscr{B}}

\newcommand{\thmIdec}{\mathscr}
\newcommand{\X}{\thmIdec X}
\newcommand{\K}{\thmIdec K}
\newcommand{\Y}{\thmIdec Y}
\newcommand{\bLowercase}{\mathfrak{b}}

\DeclareMathOperator{\coeffs}{coeffs}



%\usepackage{listings}
\definecolor{isarblue}{HTML}{006699}
\definecolor{isargreen}{HTML}{009966}
\lstdefinelanguage{isabelle}{%
    keywords=[1]{type_synonym,datatype,fun,abbreviation,definition,proof,lemma,theorem,corollary,by,have,from,with,using,lift_definition,locale,consts},
    keywordstyle=[1]\bfseries\color{isarblue},
    keywords=[2]{fixes,defines,where,assumes,shows,and,is},
    keywordstyle=[2]\bfseries\color{isargreen},
    keywords=[3]{if,then,else,case,of,SOME,let,in,O},
    keywordstyle=[3]\color{isarblue},
    commentstyle=\color{gray},
    morecomment=[l]{(*}
}
\lstset{%
  language=isabelle,
  escapeinside={&}{&},
  columns=fixed,
  extendedchars,
  basewidth={0.5em,0.45em},
  basicstyle=\ttfamily,
  mathescape,
}
%%%%%%

\begin{document}

\maketitle

\begin{abstract}
We present a universal construction of Diophantine equations with bounded complexity in Isabelle/HOL. This is a formalization of our own work in number theory~\cite{manuscript}.

Hilbert's Tenth Problem was answered negatively by Yuri Matiyasevich, who showed that there is no general algorithm to decide whether an arbitrary Diophantine equation has a solution. However, the problem remains open when generalized to the field of rational numbers, or contrarily, when restricted to Diophantine equations with bounded complexity, characterized by the number of variables $\nu$ and the degree $\delta$. If every Diophantine set can be represented within the bounds $(\nu, \delta)$, we say that this pair is \emph{universal}, and it follows that the corresponding class of equations is undecidable. In a separate mathematics article, we have determined the first non-trivial universal pair for the case of integer unknowns. 

In this paper, we contribute a formal verification of the main construction required to establish said universal pair. In doing so, we markedly extend the Isabelle AFP entry on multivariate polynomials~\cite{polynomials-afp}, formalize parts of a number theory textbook~\cite{nathanson}, and develop classical theory on Diophantine equations~\cite{MR75} in Isabelle. Additionally, our work includes metaprogramming infrastructure designed to efficiently handle complex definitions of multivariate polynomials. Our mathematical draft has been formalized while the mathematical research was ongoing, and benefitted largely from the help of the theorem prover. We reflect how the close collaboration between mathematician and computer is an uncommon but promising modus operandi.
\end{abstract}

\section{Introduction}
\subsection{Hilbert's Tenth Problem}
The tenth problem on David Hilbert's 1900 list of twenty-three mathematical problems asks for an algorithm to decide the solvability of any arbitrary Diophantine equation~\cite{hilbert-problems}. Like many problems on his list, Hilbert's Tenth Problem has provided inspiration for much mathematical research in the last century, and still does so until today. While the original problem was solved over fifty years ago~\cite{exp-diophantine}, this paper presents the formal verification of results related to a new quantitative theorem, stronger than the original solution.

A \emph{Diophantine equation} is a multivariate polynomial equation over the integers. There are countably infinitely many such equations, and they form a very general language: many difficult problems across number theory and related fields can be equivalently expressed as a single Diophantine equation~\cite{riemann-register-machine}. One often considers parametric equations of the type $P(a, z_1, \ldots, z_\nu)=0$ where $a$ is called the parameter and $z_1, \ldots, z_\nu$ are called unknowns. A \emph{Diophantine set} $A \subset \N$ is then defined to be any set of parameters such that $a \in A$ if and only if there are unknowns $z_1, \ldots, z_\nu \in \Z$ with $P(a, z_1, \ldots, z_\nu) = 0$.

The Davis-Putnam-Robinson-Matiyasevich theorem (1970)~\cite[Chapter~1]{matiyasevich-book} states that every recursively enumerable set is Diophantine; this is a deep relation between number theory and computability theory. An enumeration argument shows that, conversely, every Diophantine set is recursively enumerable. The resulting equivalence implies that there are Diophantine sets which are not decidable, and hence proves that Hilbert's Tenth Problem is undecidable.

\subsection{Universal Diophantine Equations and Pairs}
Because the set of all recursively enumerable (r.e.) sets is itself recursively enumerable, there must exist a corresponding Diophantine equation for this set, representing all r.e. sets, i.e. all Diophantine sets at once. An equation with this property is known as a \emph{universal Diophantine equation}, and an explicit example can be found in an article by Jones~\cite{9var}.

It is now natural to ask about the minimal complexity required for such universal equations. A popular measure of complexity is the number of unknowns $\nu$ required and the total degree $\delta$ of the polynomial. If all Diophantine sets can be represented by a polynomial in at most $\nu$ unknowns and in degree at most $\delta$, we call $(\nu, \delta)$ a \emph{universal pair}. There is a trade-off between the two measures: one may reduce the degree by introducing auxiliary unknowns, and vice versa eliminate unknowns at the cost of increasing the degree.

There is a subtle but interesting distinction between allowing the unknowns $z_1, \ldots, z_\nu$ of $P$ to take values in the natural numbers, including zero, versus in the integers. Universal pairs can be defined for both cases, and we shall write $(\nu, \delta)_\N$ and $(\nu, \delta)_\Z$ to make the distinction explicit. Previous work has primarily focused on universal pairs over the natural numbers.
%~\cite{9var}.
In this case, it was Matiyasevich~\cite{M77} who first announced the smallest currently known number of unknowns $\nu = 9$. Jones~\cite{9var} later published a complete proof that $(9, 1.64 \cdot 10^{45})_\N$ is universal, together with the lowest currently known degree $\delta = 4$. More precisely, the pair $(58, 4)_\N$ is universal, too.

Our work contributes the first derivation of a non-trivial universal pair over the integers. Note that any universal pair over $\N$ induces a pair over $\Z$. Every $n \in \N$ can be written as $x^2 + y^2 + z^2 + w^2$ for integers $x,y,z,w \in \Z$ using Lagrange's four square theorem. Thus, a Diophantine equation with unknowns in $\N$ can be transformed into an equation with unknowns in $\Z$ by introducing additional unknowns. However, this conversion gives bounds that are far from optimal. 

Using the four-squares substitution for $\nu = 9$ over the natural numbers would thus require 36 integer unknowns. A more efficient substitution is possible using the more sophisticated three squares theorem (cf. \cref{sec:three-squares}); this would yield universality for 27 integer unknowns. 

\subsection{First Nontrivial Bounds over \texorpdfstring{$\Z$}{the Integers}}
In this article, we formalize the main result leading to the following improvement to only eleven unknowns:

\begin{theorem}\label{thm:11-universal-pair}
    Let $(\nu, \delta)_\N$ be universal. Then
        \[ \big(11, \eta(\nu, \delta) \big)_\Z \]
    is universal where $\eta(\nu, \delta) = 226\,176 \; \delta \, (2 \delta +1)^{\nu +2} + 96 \; \delta \, (2 \delta +1)^{\nu +1} + 226\,176 \; \delta +14\,656$. \\
    In particular, 
    \[ (11,1.16258\cdot 10^{63})_\Z \qquad \text{and} \qquad (11, 9.19593 \cdot 10^{53})_\Z \]
    are universal over $\Z$.
\end{theorem}

The technique that is used to reduce the number of unknowns to eleven was developed by Sun~\cite{Sun}. However, that work does not consider the degree of the constructed polynomials.

Observe that, in principle, the second pair stated above is better than the first. We state both because the second pair depends on Jones' universal pair $(32, 12)_\N$ of which there is no published proof in the literature. The first pair depends instead on $(58, 4)_\N$.

\cref{thm:11-universal-pair} implies the following stronger response to Hilbert's Tenth Problem.
\begin{corollary}\label{cor:strong-htp}
    Consider the class of Diophantine equations with at most 11 unknowns and at most degree $1.63 \cdot 10^{63}$. Hilbert's Tenth Problem is unsolvable for this class: There is no algorithm that can determine for all equations in this class whether they have a solution or not. 
\end{corollary}

Theorem~\ref{thm:11-universal-pair} crucially rests on the following result. Its proof contains an explicit construction of the polynomial $\tilde Q$. 

\begin{theorem}\label{thm:11-unknowns}  
    Let $P \in \Z[a, z_1, \ldots, z_\nu]$. Then there is an explicit polynomial $\tilde Q \in \Z[a, z_1, \ldots, z_{11}]$ such that, for all $a \in \N$,
    \[
        \exists z_1, \ldots, z_\nu \in \Z^\nu. \; P(a, z_1, \ldots, z_\nu) = 0 \iff \exists \tilde z_1, \ldots, \tilde z_{11} \in \Z^{11}. \; \tilde Q(a, \tilde z_1, \ldots, \tilde z_{11}) = 0.
    \]
\end{theorem}

This theorem has been formalized in Isabelle/HOL and is the focus of this paper. It appears, together with Theorem~\ref{thm:11-universal-pair} and Corollary~\ref{cor:strong-htp}, in our recent manuscript~\cite{manuscript} which has been written in parallel with the development of its formal proof presented here.

\section{Related Work}
Hilbert's Tenth Problem and related statements have been formalized in a variety of different systems. Note that there are many different but equivalent options to define recursively enumerable sets, either algebraically or more by using some Turing-complete model of computation. In the following, we summarize formalizations of results related to Hilbert's Tenth Problem, highlighting the different choices authors have made.

In Lean, Carneiro~\cite{lean-formalization} has formalized that exponentiation is Diophantine, which has historically been the last step to be completed in the proof of the DPRM theorem. In Coq, Larchey-Wendling and Forster~\cite{coq-formalization} have formalized the complete DPRM theorem, elegantly using Conway's FRACTRAN language as a Turing-complete model of computation. They moreover formalize a reduction of their library of unsolvable problems, including two variants of the Halting Problem, to Hilbert's Tenth Problem. In Mizar, Pak~\cite{mizar-formalization} has formalized the DPRM theorem, defining recursively enumerable sets by Davis' normal form. In Isabelle, Bayer et~al.~\cite{isabelle-formalization} have formalized the DPRM theorem using register machines as a model of computation.

\medskip

Mathematically, we rely on techniques originally developed for working with Diophantine equations with natural number unknowns; these have been developed by Matiyasevich, Robinson~\cite{MR75} and Jones~\cite{9var}. Over the integers, Sun~\cite{Sun}
%~\cite[Theorems 3.1, 4.1 \& 4.2 and 1.1]{Sun}
first announced many results. In particular, he proved that any Diophantine equation can be reduced to an equivalent one with only eleven integer unknowns. 

Unfortunately, the published literature often lacked sufficient detail to formalize, and sometimes even to understand, the existing proofs. As such, our detailed mathematical manuscript now reproduces many of the results and proofs cited above. The improved clarity in exposition is a notable product of this formalization. 

\section{Preliminaries}

\subsection{Multivariate Polynomials in Isabelle}
Throughout our work, we rely on the formalization of executable multivariate polynomials by Sternagel et al.~\cite{polynomials-afp}. The type of multivariate polynomials is defined as
\begin{align}\label{eq:typedef-mpoly}
    \alpha\; \texttt{mpoly} := (\N \Rightarrow_0 \N) \Rightarrow_0 \alpha
\end{align}
for any type $\alpha$ with a zero. Most commonly we will take $\alpha = \Z$ or a unital, commutative ring. We use $\beta \Rightarrow_0 \alpha$ to denote a \emph{finite assignment}, i.e. a function from $\beta$ to $\alpha$ that takes only finitely many nonzero values. Each finite assignment $m : \N \Rightarrow_0 \N$ defines an abstract multivariate monomial $\prod_i x_i^{m(i)}$, giving the exponents of each $x_i$, for $i\in \N$. A multivariate polynomial then assigns a coefficient from $\alpha$ to each monomial. Of course, this type inherits much algebraic structure from the underlying coefficient type: if $\alpha$ is a ring or an integral domain, then so is $\alpha$ \texttt{mpoly}. With this formalization, we extend the existing library by well over 2,000 lines of proof.

A key first addition to the existing library is an explicit expansion, for arbitrary $\nu \in \N$, of a general multivariate polynomial $P \in \alpha[z_0, \ldots, z_\nu]$ as
\begin{equation}\label{eq:multivariate-expansion}
    P(z_0, \ldots, z_\nu) = \sum_{i_0, \ldots, i_\nu \, \in \, \N} c_{i_0, \ldots, i_\nu} \prod_{s=0}^\nu z_s^{i_s}.
\end{equation}
This expression is a polynomial as long as only finitely many coefficients $c_{i_0, \ldots, i_\nu}$ are non-zero. Below, we abbreviate the indices with a single multi-index $\mathbf{i} = (i_0, \ldots, i_\nu) \in \N^{\nu+1}$, defining $\|\mathbf{i}\| := \sum_{s=0}^\nu i_s$.

With $\delta := \deg P$ and the constraint $\|\mathbf{i}\| \leq \delta$, we can reformulate the same expression as a finite sum, which is much better behaved for formal proofs in Isabelle.
\begin{equation}\label{eq:multivariate-expansion-bounded}
    P(z_0, \ldots, z_\nu) = \sum_{\substack{\mathbf{i} \in \N^{\nu+1} \\ \|\mathbf{i}\| \le \delta}} c_{i_0, \ldots, i_\nu} \prod_{s=0}^\nu z_s^{i_s},
\end{equation}

Although we will not provide typeset proofs of equations~\eqref{eq:multivariate-expansion} and~\eqref{eq:multivariate-expansion-bounded}, neither equation follows trivially from the definition~\eqref{eq:typedef-mpoly} of the type $\alpha$ \texttt{mpoly}. In Isabelle, we have formalized equation~\eqref{eq:multivariate-expansion-bounded} as follows. The term \texttt{Const C} represents a constant polynomial with value $C$, and the term \texttt{Var s} represents the linear polynomial $z_s$. The term \texttt{max\_vars P} gives the highest variable index that appears in $P$, denoted $\nu$ above.

\needspace{4\baselineskip}
\begin{lstlisting}
lemma mpoly_multivariate_expansion:
  fixes P :: "$\alpha$ mpoly"
  shows "P = $\sum$i | length i = max_vars P + 1 $\wedge$ ||i|| $\leq$ total_degree P.
          Const (coeff P i) * ($\prod$s = 0..max_vars P. (Var s) ^ (i ! s)))"
\end{lstlisting}

\begin{remark}
    In the interest of legibility, throughout the paper, we are presenting Isabelle-style pseudo-code in which we have dropped some notation that is syntactically but not semantically necessary.
    % This includes extra parentheses and the morphisms associated to an Isabelle \texttt{typedef}, as well as conversions between the two (equivalent) types $\tau$ \texttt{list} of finite lists and \texttt{nat} $\Rightarrow_0 \tau$ of finite assignments indexed by a natural number.
    Moreover, by writing \texttt{||i||} above we mean the library function \texttt{sum\_list i}. We write $\alpha$ for the free type variable \texttt{'a} representing the ring of coefficients. We assume $\alpha$~\texttt{::} \texttt{comm\_ring\_1}, a commutative and unital ring, throughout.
\end{remark}


In the same way, we developed other basic notation and tools to efficiently reason about the type $\alpha$ \texttt{mpoly}. For example, we formalize that if one type injects into a larger type like $\Z \hookrightarrow \mathbb{Q}$, then the same holds for $\Z$ \texttt{mpoly} $\hookrightarrow \mathbb{Q}$ \texttt{mpoly}. We also prove their correct behavior with respect to standard operations on polynomials: addition, multiplication, exponentiation et cetera.

\paragraph*{Substitution}
One property that has been vital for the iterative definition of universal Diophantine constructions is the substitution of polynomials into polynomials. The existing AFP entry provides a morphism \texttt{reduce\_nested\_mpoly ::} \texttt($\alpha$ \texttt{mpoly) mpoly} $\Rightarrow$ $\alpha$ \texttt{mpoly}. However, we instead chose to formalize substitutions more directly and more explicitly, adapted to the constructive proof presented in the next section.

\begin{lstlisting}
definition poly_subst_monom :: "(nat$\Rightarrow$$\alpha$ mpoly) $\Rightarrow$ (nat$\Rightarrow_0$nat) $\Rightarrow$ $\alpha$ mpoly"
  where "poly_subst_monom f m = $\prod$s. (f s) ^ (m s)"

definition poly_subst :: "(nat$\Rightarrow$$\alpha$ mpoly) $\Rightarrow$ $\alpha$ mpoly $\Rightarrow$ $\alpha$ mpoly"
  where "poly_subst f P = $\sum$m. (Const (coeff P m)) * poly_subst_monom f m"
\end{lstlisting}
This formalizes an expansion of the polynomial $P$ in analogy to the above lemma, except that each \texttt{Var s} has been replaced by corresponding substitution $f_s$, which is now allowed to be any arbitrary multivariate polynomial. 

In standard mathematical notation, given a polynomial $P \in \Z[z_0, \ldots, z_\nu]$, we would write this substitution as a list of $(\nu+1)$-many polynomials $f = (f_0, \ldots, f_\nu) \in \Z[z_0, \ldots, z_\nu]^{\nu+1}$. One often requires the corresponding correctness statement: the order of evaluation does not matter. That means, given a vector of integers $\mathbf{z} = (z_0, \ldots, z_\nu) \in \Z^{\nu+1}$, we need to prove that
\begin{equation}\label{eq:substitution-evaluation}
    P(f)(\mathbf{z}) = P(f(\mathbf{z})),
\end{equation}
where $f(\mathbf{z}) = (f_0(\mathbf{z}), \ldots, f_\nu(\mathbf{z})) \in \Z^{\nu+1}$ is the evaluation of each substitution polynomial at $\mathbf{z}$. In Isabelle, the evaluation morphism is formalized as \texttt{insertion z P} for an assignment of the type \texttt{z :\!\!\! nat $\Rightarrow \alpha$} and a multivariate polynomial \texttt{P :\! $\alpha$ mpoly}. Then, equation~\eqref{eq:substitution-evaluation} reads:
\begin{lstlisting}
lemma insertion_poly_subst:
  "insertion z (poly_subst f P) = insertion ((insertion z) $\circ$ f) P"
\end{lstlisting}


\paragraph*{Number of Variables and Total Degree}
To explicitly formalize universal pairs, we need to be able to calculate the degree and number of variables of a multivariate polynomial in Isabelle. The AFP entry~\cite{polynomials-afp} defines both \texttt{vars} and \texttt{total\_degree}, to which we have added the definition \texttt{max\_vars}.
\begin{lstlisting}
lift_definition vars :: "$\alpha$ mpoly $\Rightarrow$ nat set"
  is "$\lambda$P. $\bigcup$ (keys ` keys P)"

definition max_vars :: "$\alpha$ mpoly $\Rightarrow$ nat"
  where "max_vars P $\equiv$ Max (insert 0 (vars P))"

lift_definition total_degree :: "$\alpha$ mpoly $\Rightarrow$ nat"
  is "$\lambda$P. Max (insert 0 (($\lambda$m. sum (lookup m) (keys m)) ` keys P))"
\end{lstlisting}
The morphism \texttt{keys} is defined to be the set of all elements of the domain which are mapped to a non-zero value. In particular, \texttt{keys P} is the set of all monomials which appear in $P$ with non-zero coefficient. The morphism \texttt{lookup} denotes function evaluation for finite assignments. Recall also that \texttt{sum f S} $= \sum_{s\in S} f(s)$.

In Isabelle, the \texttt{Max} of a set is not defined for empty sets. Because any constant polynomial has \texttt{vars (Const c)} $= \emptyset$, inserting the number $0$ in the definition of \texttt{max\_vars} ensures that the set is always non-empty. When constructing polynomials we always make sure to not skip any variable indices. Moreover, because we work with parametric polynomials, the zeroth variable is always reserved for the parameter $a$. Accordingly, the above definition has the correct semantics for this project.

The definition of \texttt{total\_degree} works in analogy to \texttt{max\_vars}. Because the zero polynomial has \texttt{keys (Const 0)} $= \emptyset$, we again include \texttt{insert 0}. However, this formalizes the degree of the zero polynomial as~$0$. A consequence of this definition is, e.g., that the assumptions $P \neq 0$ and $Q \neq 0$ are now required for the lemma \texttt{total\_degree (P * Q) = total\_degree P + total\_degree Q}.

\medskip

%In order to actually calculate the number of variables used in Theorem~\ref{thm:11-unknowns}, 
We formalize appropriate simplification and correctness properties of these definitions for all other functions and morphisms which we use to construct polynomials. This includes exponentiation, finite sums and products, type casting, and substitutions. Because we are concerned with upper bounds of the degree only, it suffices to establish inequalities which require fewer assumptions. Indeed, \texttt{total\_degree (P * Q) $\leq$ total\_degree P + total\_degree Q} is true without any assumptions on $P$ or $Q$.

With respect to the morphism \texttt{poly\_subst} introduced above, for example, we formalize the following relative upper bounds. Note that \texttt{degree P s} calculates the degree of $P$ in the variable $z_s$ only.
\begin{lstlisting}
lemma vars_poly_subst:
  "vars (poly_subst f P) $\subseteq$ $\bigcup$s $\in$ vars P. vars (f s)"
  (* Proof omitted *)
\end{lstlisting}
\needspace{4\baselineskip}
\begin{lstlisting}
lemma total_degree_poly_subst: 
  "total_degree (poly_subst f P) 
      $\leq$ $\sum$s $\in$ vars P. degree P s * total_degree (f s)"
  (* Proof omitted *)
\end{lstlisting}



\paragraph*{Custom Metaprogramming Infrastructure}
Recall that \cref{thm:11-unknowns} asserts that any Diophantine equation can be expressed as a polynomial $\tilde Q$ with only eleven integer unknowns. Consequently, the natural choice for the type of $\tilde Q$ in Isabelle is \texttt{int mpoly}. Since $\tilde Q$ is constructed by substituting various previously defined polynomials into each other, one might expect that these intermediate polynomials also have the type \texttt{int mpoly}.

However, this approach presents notable difficulties when formalizing the \cref{thm:coding,thm:lucas} which will be presented in the following section. Instead of using meaningful variable names like $a, f, g$, one must work with \texttt{Var 0}, \texttt{Var 1}, \texttt{Var 2} and frequently refer to the insertion map explicitly. More critically, substituting polynomials into one another requires care to avoid collisions of variable indices. These challenges cannot be fully mitigated through new notation alone.

To address these issues, this project developed a custom metaprogramming framework that effectively bridges Isabelle terms and the type \texttt{int mpoly}. This allows our formalization to benefit from the best of both worlds, significantly simplifying our proofs. 

Using this framework, a polynomial $P$ is initially formalized as a standard Isabelle term of type \lstinline{P :: int => int => ... => int}. The command \texttt{poly\_extract P}
%, implemented by Anton Danilkin,
then generates a new term \texttt{P\_poly} of type \texttt{int mpoly} and automatically proves a lemma \texttt{P\_correct}. This lemma guarantees that evaluating \texttt{P\_poly} yields the same result as evaluating the original function \texttt{P}. The extraction succeeds if \texttt{P} is a valid polynomial: constructed from integer constants and variables, addition, multiplication, and any polynomial term \texttt{C} for which a corresponding correctness lemma \texttt{C\_correct} exists.

\subsection{The Three Squares Theorem} \label{sec:three-squares}
The three squares theorem states that for every $n \in \N$ there exist integers $x, y, z \in \Z$ such that $n = x^2 + y^2 + z^2 + z$. It is a consequence of the classical number theory developed in Nathanson's textbook~\cite[Lemma~1.9]{nathanson}. Danilkin and Chevalier formalized the three squares theorem as well as its number theory prerequisites in Isabelle~\cite{three-squares-afp}. Using this result will constitute the last step in our proof. It serves to ultimately replace the single remaining non-negative integer $n \geq 0$ in Lemma~\ref{lem:relation-combining} below by three true integer unknowns.

\subsection{Relation Combining}
The idea of relation-combining is the following: imagine we want to check the two conditions $x > 0, y \ne 0$ with Diophantine polynomials in integer variables. A naive method would use the three-squares formula for each condition, i.e. the conditions are true if
\begin{align*}
x &= a^2+b^2+c^2+c+1 \\
y^2 &= d^2+e^2+f^2+f+1
\end{align*}
for some $a,b,c,d,e,f \in \Z$.

However, there is a more efficient method because the two conditions together are equivalent to the single condition
\[
xy^2 = a^2+b^2+c^2+c+1
\;,
\]
using only $a, b, c \in \Z$, cutting the number of integer variables required in half.

This technique has been mastered by Matiyasevich and Robinson with the following result. It combines a divisibility condition, an inequality, and checking whether integers are perfect squares into the existence of only one (!) non-negative integer.

\begin{lemma}[Matiyasevich-Robinson~\cite{MR75}]\label{lem:relation-combining}
    For every $q>0$ there is a polynomial $M_q$ with the property that for all integers $A_1,\dots,A_q,R,S,T$ such that $S\neq 0$, the following are equivalent:
    \begin{enumerate}
        \item $S|T$, $R>0$ and $A_1,\dots,A_q$ are squares,
        \item $\exists n \ge 0 : \, M_q(A_1,\dots,A_q,S,T,R,n)=0$.
    \end{enumerate}
    Specifically, $M_q$ can be given as	
    \begin{equation} \label{eq:matiyasevich-polynomial} \begin{split}
        &M_q(A_1,\dots,A_q,S,T,R,n)\\
        &=\prod_{\varepsilon_1,\dots,\varepsilon_q\in\{\pm1\}^q}\left(S^2n + T^2 - S^2(2R-1)\left(T^2 + X_q^q +\sum_{j=1}^q \varepsilon_j \sqrt{A_j}X_q^{j-1}\right) \right) \\
    \end{split} \end{equation}
    where $X_q:=1+\sum_{j=1}^{q}A_j^2$.
\end{lemma}

The short-hand notation in equation~\eqref{eq:matiyasevich-polynomial} defines a polynomial because $M_q$ is even in $\sqrt{A_j}$. In our case, we do not require the general form of this result and instead only its specialization to $q = 3$. Using the automation built into \texttt{poly\_extract} as discussed above, it moreover suffices to simply spell out the explicit function $\texttt{M3} : \Z^7 \rightarrow \Z$. Lemma~\ref{lem:relation-combining} is then formalized as follows.

\needspace{6\baselineskip}
\begin{lstlisting}
lemma relation_combining:
  assumes "S $\neq$ 0"
  shows "(S dvd T $\wedge$ R > 0 $\wedge$ is_square A${}_1$ $\wedge$ is_square A${}_2$ $\wedge$ is_square A${}_3$)
         = ($\exists$n$\geq$0. M3 A${}_1$ A${}_2$ A${}_3$ S T R n = 0)"
\end{lstlisting}

\section{Structure of the Proof}
\cref{thm:11-universal-pair} is based on three technical results that capture the core ideas behind universal Diophantine constructions: Coding, Lucas Sequences, and relation combining techniques. The dependencies between these results are illustrated in \cref{fig:proof-structure}.

The central challenge in obtaining a universal pair is reducing a general Diophantine equation --- potentially involving many unknowns --- to an equivalent equation with a fixed number of unknowns. This reduction is carried out in the aforementioned steps. In the process, the original equation is first transformed into intermediate expressions, which are then shown to be equivalent to a fixed Diophantine polynomial.

Crucially, all operations appearing in the intermediate expressions must be Diophantine, that is, expressible using polynomials. Many fundamental mathematical relations admit such representations~\cite{matiyasevich-book}, including divisibility, congruence modulo an integer, and inequalities. Relations involving exponentiation, binomial coefficients and factorials also admit Diophantine representations. However, when deriving universal pairs, it is not enough to establish mere existence; the representation must also be efficient, meaning that it should involve a minimal number of unknowns and remain of low degree.

\begin{figure}[omitted: a diagram of this paper's results, not reproduced here]\label{fig:proof-structure}\end{figure}

\section{Coding}
Coding is the key idea that translates a given Diophantine polynomial $P \in \Z[z_1, \ldots, z_\nu]$ with a large number of unknowns into a Diophantine relation of bounded complexity, i.e. only requiring a fixed number of unknowns.
%We first give the definition of a code and then explain how to achieve the reduction in unknowns. 

\begin{definition}[Code] \label{def:code}
Fix a base $\B = 2^k$ for some integer $k > 1$. Let $z_1, \ldots, z_\nu$ be given and $z_i < \B$ for all $i=1, \ldots, \nu$. We say that $c$ is a code for the tuple $(z_1, \ldots, z_\nu)$ with respect to a tuple $(n_1, \ldots, n_\nu)$ of strictly increasing numbers if $c = \sum_{i=1}^\nu z_i \B^{n_i}$.
\end{definition}

If the $z_i$ are the given unknowns of the polynomial $P$ then $c$ consolidates them in a single large number. It is possible to define a number $\alpha_P$ such that the product $K = (1 + c)^{\deg P} \cdot \alpha_P$ contains the value $P(z_1, \ldots, z_\nu)$ in a specific location of its base $\B$ expansion.\footnote{More precisely, such $\alpha_P$ exists if $P$ only contains nonnegative coefficients; the statement can be generalized to arbitrary $P$ by adding an additional correcting term (cf. \cite[Lemma 1.8]{manuscript} where $\alpha_P$ appears as $\coeffs(\B)$).} One can therefore translate the statement $P(z_1, \ldots, z_\nu) = 0$ into a statement about the base $\B$ expansion of $K$. Ensuring that a certain base $\B$ digit of $K$ is zero can then be encoded by counting carries during binary addition. In turn, this condition can be translated into a divisibility condition on a binomial coefficient~\cite[Lemmas~1.3 and~1.4]{manuscript}. 

To formalize codes and work with natural numbers in base $\B$ we make use of the AFP entry on digit expansions~\cite{afp-digit-expansions}. We also provide a formal definition of the factorial for tuples of natural numbers and formalize the multinomial theorem --- these are required for working with $\alpha_P$.

\paragraph*{The Coding Theorem}
To combine the various translation steps into a single theorem a comprehensive construction is required. One iteratively defines the auxiliary variables $\mathcal{L}, r, \beta, \alpha, \gamma$ and constructs polynomials $\mathfrak{b}, \mathscr{B}, M, N_0, N_1, N, \mathscr{T}, \Y \in \Z[a, f]$ as well as $c, \mathscr{K}, \mathscr{S}, \mathscr{R}, \X \in \Z[a, f, g]$. All quantities depend on a given polynomial $P \in \Z[a, z_1, \ldots, z_\nu]$ with parameter $a$ and $\nu$ unknowns. In particular, $\K$ corresponds to $K$ given above. We skip the precise definitions here for the sake of brevity (cf. \cite[Definition 2.1]{manuscript}). 

% CHANGED by MD
%The following statement of the Coding Theorem uses the notation $[j, k) := \{ n \in \N : j \leq n < k \}$.
The following statement uses the notation $[j, k) := \{ n \in \N : j \leq n < k \}$.

\begin{theorem}[Coding Theorem~{\cite[Theorem 1]{manuscript}}] \label{thm:coding}
For $a \ge 0$ and $f \neq 0$ such that $\bLowercase := \bLowercase(a, f)$ is a power of two the following are equivalent:
\begin{equation} \label{eq:polynomial-zero}
	\exists (z_1,\ldots,z_\nu)\in [0, \bLowercase)^\nu \colon 
    P(a, z_1, \ldots, z_\nu) = 0 
\end{equation}
and
\begin{equation} \label{eq:divisibility-binomial}
	\exists g \in [\bLowercase, \gamma \bLowercase^\alpha) \colon  \Y(a, f) \, |\, \binom{2 \X(a, f, g)}{\X(a, f, g)} \;.
\end{equation}
Moreover, \eqref{eq:divisibility-binomial} $\implies $ \eqref{eq:polynomial-zero} already holds when $g \in [0, 2\gamma \bLowercase^\alpha)$.
\end{theorem}

Recall that divisibility, exponentiation and binomial coefficients are Diophantine. We have therefore reduced the original polynomial with $\nu$ unknowns to a Diophantine relation in the parameter $a$ and only two unknowns $f, g$.

Our formalization of this theorem employs two locales. The first locale \texttt{coding\_variables} fixes the polynomial $P$ along with two integers $a$ and $f$. Within this locale, the polynomials $\bLowercase, \X,\Y$, as well as the auxiliary variables mentioned above, are introduced. The second locale, \texttt{coding\_theorem}, incorporates the assumptions of \cref{thm:coding}, which are also used in several lemmas leading up to its proof. This structured approach proves advantageous in later stages of the formalization, ensuring that all coding-related concepts are accessed through a unified interface. 
\begin{lstlisting}
locale coding_theorem = coding_variables + 
  assumes a_nonneg: "a $\geq$ 0"
      and f_pos: "f > 0"
      and b_power2: "is_power2 $\bLowercase$"
      and $\delta$_pos: "$\delta$ > 0"
\end{lstlisting}

Within this locale we formalize two separate statements corresponding to the two directions of the equivalence in \cref{thm:coding}. Here we present the statement of the reverse direction. Because $a$ and $f$ are fixed in the locale, $\X$ becomes a univariate polynomial in $g$ in the below code.
\begin{lstlisting}
theorem coding_theorem_reverse':
  assumes "$\exists$g. 0 $\leq$ g 
            $\land$ g < 2*(int $\gamma$)*$\bLowercase$^$\alpha$ 
            $\land$ $\Y$ dvd (2 * nat ($\X$ g) choose nat ($\X$ g))"
    shows "$\exists$z. z 0 = a 
            $\land$ ($\forall$i. 0 $\leq$ z i $\land$ z i < $\bLowercase$) 
            $\land$ insertion z P = 0"
\end{lstlisting}

A significant portion of the proof of this statement involves establishing bounds between the various variables and polynomials. The inequalities written in our mathematical draft were often not accurate or they were missing crucial assumptions. It was very convenient to then work out the details in Isabelle which tracked precisely which bounds had already been established and under which side conditions --- for this type of mathematics, the proof \emph{assistant} lived up to its name! We made heavy use of literal facts (writing \texttt{\guilsinglleft b > 0\guilsinglright} instead of using an identifier like \texttt{b\_gt\_0}) which additionally improved the presentation of the proof. 

\section{Lucas Sequences and the Pell Equation}
The Coding Theorem introduces an exponential relation --- namely that the base $\bLowercase$ be a power of two --- and a binomial coefficient. Both need to be converted into an equivalent polynomial condition. To do so efficiently, we rely on the solutions of a Pell equation, also known as Lucas sequences. 

\begin{definition}[Lucas sequence]
A \emph{Lucas sequence with parameter $A\in\Z$} is a sequence $(x_n)_{n\in \Z}$ that satisfies the recursive relation 
\begin{equation}\label{eq:lucas-general-recurrence}
    x_{n+1}=Ax_n-x_{n-1}.
\end{equation}
Each such sequence is defined uniquely by two initial values, and we define two particular Lucas sequences with initial values $\psi_0(A)=0, \psi_1(A)=1$ as well as $\chi_0(A)=2, \chi_1(A)=A$.
\end{definition}
In Isabelle, we first define these Lucas sequences for $n \in \N$.
\begin{lstlisting}
fun $\psi$ :: "int $\Rightarrow$ nat $\Rightarrow$ int" where
  "$\psi$ A 0 = 0" | "$\psi$ A (Suc 0) = 1" |
  "$\psi$ A (Suc (Suc n)) = A * ($\psi$ A (Suc n)) - ($\psi$ A n)"

fun $\chi$ :: "int $\Rightarrow$ nat $\Rightarrow$ int" where
  "$\chi$ A 0 = 2" | "$\chi$ A (Suc 0) = A" |
  "$\chi$ A (Suc (Suc n)) = A * ($\chi$ A (Suc n)) - ($\chi$ A n)"
\end{lstlisting}
We then extend this recursive definition to all $n \in \Z$ by using the symmetries of these sequences, given by $\psi_{-n}(A) = -\psi_n(A)$ and $\chi_{-n}(A) = \chi_n(A)$. With these properties, it follows directly that Equation~\eqref{eq:lucas-general-recurrence} holds for the entire integer-indexed sequence.

\begin{lstlisting}
definition $\psi$_int :: "int $\Rightarrow$ int $\Rightarrow$ int"
  where "$\psi$_int A n = (if n $\geq$ 0 then 1 else -1) * $\psi$ A (nat (abs n))"

definition $\chi$_int :: "int $\Rightarrow$ int $\Rightarrow$ int"
  where "$\chi$_int A n = $\chi$ A (nat (abs n))"
\end{lstlisting}

These two second-order recurrences grow exponentially: for example, if $n \ge 2$ and $A>1$, we have $(A-1)^n < \psi_{n+1}(A) < A^n$. This is the property we will use to encode that $\bLowercase$ is a power of two. Importantly, these recurrences are Diophantine, in the sense that they are solutions of the following polynomial equation.

\begin{lemma}[Pell's equation]
\label{lem:PellEquation}
If $d=A^2-4$, the Pell equation
\begin{equation}
    X^2-dY^2=4
\end{equation}
has the property that its solutions are exactly the pairs $(X,Y)=(\chi_n(A),\psi_n(A))$ for $n\in\Z$.
\end{lemma}
The following formal statements are sufficient for our purposes. 

\begin{lstlisting}
lemma lucas_solves_pell: "(A${}^2$-4)*($\psi$_int A m)${}^2$ + 4 = ($\chi$_int A m)${}^2$"
lemma pell_yields_lucas: "($\exists$k. (A${}^2$-4)*Y${}^2$ + 4 = k${}^2$) = ($\exists$m. Y = $\psi$_int A m)"
\end{lstlisting}
In the course of the proofs of these statements, totaling about 2,000 lines, we formalize many auxiliary properties of Lucas sequences. This includes both a wealth of elementary facts (cf.~\cite[Lemma~3.4]{manuscript}) and less trivial properties (cf.~\cite[Lemmas~3.7--3.12]{manuscript}). All of the latter had previously been stated or sketched in the literature, but this formalization, and our associated manuscript, contain their complete proofs for the first time.

\paragraph*{A Jungle of Variables to Build a Bridge}

The specfic relations which will represent the exponential conditions from the previous section (power of two and binomial coefficient) rely on a dozen auxiliary definitions. This is not the only list of this length; also the Coding Theorem~\ref{thm:coding} as well as the final construction in Theorem~\ref{thm:9-unknowns-ZN} depend on comparable jungles of variables. We give Definition~\ref{def:VariableDefinition} here as an example that illustrates this complexity.

It is only partially possible to assign high-level meaning to the variable definitions below. Note that if $D$ is a square, then \cref{eq:Def_D} corresponds to a Pell equation with $d = A^2 - 4$. Moreover, the lemma \texttt{pell\_yields\_lucas} can be applied to relation~\eqref{eq:UVK} to deduce that $J$ is the value of a Lucas sequence with parameter $U^2V$. Many of the remaining variables serve primarily as scaffolding, ensuring the necessary conditions to apply the appropriate number-theoretic lemmas.

\begin{definition}
\label{def:VariableDefinition}
Given variables $X,Y,b,g,h,k,l,w,x,y$ we define the following variables:
\begin{subequations}\label{eqs:VariableDefinition}
\begin{align}
U &:= 2lXY \label{eq:Def_U} \\
V &:= 4gwY \label{eq:Def_V} \\
A &:= U(V+1) \label{eq:Def_A} \\
B &:= 2X+1 \label{eq:Def_B} \\
C &:= B+(A-2)h \label{eq:Def_C} \\
D &:= (A^2-4)C^2+4 \label{eq:Def_D} \\
E &:= C^2Dx \label{eq:Def_E} \\
F &:= 4(A^2-4)E^2+1 \label{eq:Def_F} \\
G &:= 1+CDF-2(A+2)(A-2)^2E^2 \label{eq:Def_G} \\
H &:= C+BF+(2y-1)CF \label{eq:Def_H} \\
I &:= (G^2-1)H^2+1 \label{eq:Def_I} \\
J &:= X+1+k(U^2V-2) \label{eq:Def_J}
\end{align}
\end{subequations}
Denote the set of squares by $\square$. We define the following relations on the just-constructed variables:
\begin{subequations}\label{eqs:thm2-properties}
\begin{align}
&DFI \in \square \label{eq:DFI} \\
&(U^{4}V^2-4)J^2+4\in\square \label{eq:UVK} \\
&(2A-5) \,\big|\, (3bwC-2(b^2w^2-1)) \label{eq:pAppWW_new}\\
& \left(\frac{C}{J}-lY\right)^2 < \frac{1}{16g^2} \label{eq:g}
\;.
\end{align}
\end{subequations}
\end{definition}
Notice that this list contains exactly those relations which can be encoded by the Matiyasevich polynomial in Lemma~\ref{lem:relation-combining}, namely multiple perfect squares, one divisibility condition, and one inequality. This is what earns our most technical theorem its name, for it bridges the gap between the Coding Theorem~\ref{thm:coding} and the Matiyasevich polynomial.

\begin{theorem}[Bridge Theorem~{\cite[Theorem 2]{manuscript}}] \label{thm:lucas}
Given integers $b\ge 0$, $X\geq 3b$, $Y\geq \max \{ b,2^8 \}$ and $g\ge 1$, the following are equivalent:
\begin{enumerate}
\item \label{proofitem:Binom} $b$ is a power of 2 and $Y | \binom{2 X}{X}$;
\item \label{proofitem:Dioph2} $\exists h,k,l,w,x,y\geq b$ such that the Diophantine conditions \eqref{eq:DFI}, \eqref{eq:UVK}, \eqref{eq:pAppWW_new}, \eqref{eq:g} hold.
\end{enumerate}
Moreover, the implication \eqref{proofitem:Binom} $\implies$ \eqref{proofitem:Dioph2} holds already if $X\ge b$ (rather than $X\ge 3b$).
\end{theorem}

\begin{remark}
    The slightly simplified version of Theorem~\ref{thm:lucas} stated here illustrates the role of this theorem for the big picture. In full detail (cf.~\cite[Theorem~2]{manuscript}), it establishes the equivalence of three statements, where the third statement is a slightly modified version of \cref{proofitem:Dioph2}. 
\end{remark}

% \begin{remark}
%     The attentive reader may have noticed that we formalize Lemma~\ref{lem:relation-combining} to encode three perfect squares, but the conditions \eqref{eq:DFI}--\eqref{eq:g} only give two. Once we combine all previous results, we will add the third condition $\bLowercase \in \square$, simply because this will ensure $b \geq 0$ at no cost of additional variables.
% \end{remark}

This statement is formalized in Isabelle as presented below. The lists of equations~\eqref{eqs:VariableDefinition} and \eqref{eqs:thm2-properties} are implemented in the locale \texttt{bridge\_variables}. Notice that all relations are implicitly functions of the variables $X,Y,b,g,h,k,l,w,x,y$ listed in Definition~\ref{def:VariableDefinition}. To ease legibility below, we color these variables in gray whenever they are arguments to another function.

\begin{lstlisting}[escapechar={|}]
definition (in bridge_variables) bridge_relations where
  "bridge_relations k l w h x y g Y X =
      is_square (D |\color{gray}l w h g Y X| * F |\color{gray}l w h x g Y X| * I |\color{gray}l w h x y g Y X|)
      $\wedge$ is_square ((U |\color{gray}l X Y|^4 * V |\color{gray}w g Y|^2 - 4) * J |\color{gray}k l w g Y X|^2 + 4)
      $\wedge$ S |\color{gray}l w g Y X| dvd T |\color{gray}l w h g Y X b|
      $\wedge$ (4*g*C |\color{gray}l w h g Y X| - 4*g*l*Y*J |\color{gray}k l w g Y X|)^2 < (J |\color{gray}k l w g Y X|)^2"

theorem (in bridge_variables) bridge_theorem_simplified:
  fixes b Y X g :: int
  assumes "b $\geq$ 0" and "Y $\geq$ b" and "Y $\geq$ 2^8" and "X $\geq$ 3*b" and "g $\geq$ 1"
  shows "(is_power2 b $\wedge$ Y dvd (2 * X choose X))
            = ($\exists$h k l w x y :: int. bridge_relations k l w h x y g Y X
                $\wedge$ (h$\geq$b) $\wedge$ (k$\geq$b) $\wedge$ (l$\geq$b) $\wedge$ (w$\geq$b) $\wedge$ (x$\geq$b))"
\end{lstlisting}


%%%%%%%%%%%%%%%%%%%%%%%%%


\section{Eleven Unknowns in \texorpdfstring{$\Z$}{Z}}

In the previous sections, we have detailed multiple translation theorems. \Cref{thm:coding} translates between an arbitrary polynomial $P$ having a (parametric) solution and a Diophantine predicate involving exponentiation and a binomial coefficient. \Cref{thm:lucas} translates exponentiation and the binomial coefficient into more elementary Diophantine relations. Finally, \Cref{lem:relation-combining} allows us to efficiently represent precisely these relations through a single polynomial. 

\paragraph*{An Intermediate Theorem}
Before we can prove our two main Theorems~\ref{thm:11-universal-pair} and~\ref{thm:11-unknowns} stated in the introduction, we first combine these three statements into an intermediate result. Let a polynomial $P$ with natural number unknowns be given. Our goal is to construct a polynomial $Q$ with a fixed number of unknowns which should have a solution whenever $P$ has a solution. 

We do not give the elaborate definition of $Q$ here; it specifies how the various polynomials defined previously are substituted into each other~\cite[Definition 5.2]{manuscript}. In our construction, $Q$ has nine unknowns, eight of which are integers and the remaining one is a natural number. One can then prove the following theorem:

\begin{theorem}[9 Unknowns in $\Z$ and $\N$~{\cite[Theorem 3]{manuscript}}] \label{thm:9-unknowns-ZN}
For all $a \in \N$ the following are equivalent: 
\begin{gather} \label{eq:P-solution}
\exists z_1, \ldots, z_\nu \in \N \colon 
    P(a, z_1, \ldots, z_\nu) = 0 \\
\exists f, g, h, k, l, w, x, y \in \Z, n \in \N \colon 
    Q(a, f, g, h, k, l, w, x, y, n) = 0 \;. 
    \label{eq:Q-solution}
\end{gather}
\end{theorem}

The proof of this theorem proceeds by combining \cref{thm:coding,thm:lucas,lem:relation-combining}. We structure the formalization of \cref{thm:9-unknowns-ZN} into two Isabelle theories: the first one contains all definitions leading up to $Q$ and the second one implements the proof of the theorem itself. When defining $Q$ we make frequent use of the \texttt{poly\_extract} command. We introduce custom syntax to avoid having to spell out all nine unknowns repeatedly. 

We formalize two versions of the statement of \cref{thm:9-unknowns-ZN}. In the first version, $Q$ is simply an Isabelle term. The second version presented below uses $Q$ of type \texttt{int mpoly}. Note the use of \texttt{is\_nonnegative} (defined in the obvious way) which we use because $P$ has natural number unknowns. The syntax \lstinline{z(0 := a)} is used to substitute the parameter $a$ into the first variable of \texttt{P}. 
\begin{lstlisting}
theorem nine_unknowns_over_Z_and_N: 
  fixes P :: "int mpoly"
  shows "($\exists$z :: nat $\Rightarrow$ int. is_nonnegative z $\land$ 
              insertion (z(0 := int a)) P = 0) 
       = ($\exists$ f g h k l w x y n. n $\geq$ 0 $\land$
              insertion ((!) [int a, f, g, h, k, l, w, x, y, n])
              (combined_variables.Q_poly P) = 0)"
\end{lstlisting}

\paragraph*{Transforming the Last Natural Number Unknown}
The statement of \cref{thm:9-unknowns-ZN} is close to that of \cref{thm:11-unknowns}, the remaining step being that $Q$ has one remaining natural number unknown which needs to be converted into integer unknowns. One can introduce additional unknowns $n_1, n_2, n_3 \in \Z$ and replace $n$ with $n_1^2 + n_2^2 + n_3^2 + n_3$ to obtain a polynomial $\tilde Q$:
\[ \tilde Q (a, f, g, h, k, l, w, x, y, n_1, n_2, n_3) = Q(a, f, g, h, k, l, w, x, y, n_1^2 + n_2^2 + n_3^2 + n_3) \]
We can use this to prove \cref{thm:11-unknowns}.
\begin{proof}[Proof of \cref{thm:11-unknowns}]
    By the Three Squares Theorem, any $n \geq 0$ can be represented as $n_1^2 + n_2^2 + n_3^2 + n_3$. Therefore, $\tilde Q$ has a solution whenever $n \geq 0$ and $Q$ has a solution, i.e. when \eqref{eq:Q-solution} holds. Using \cref{thm:9-unknowns-ZN}, we can thus deduce that $P$ has a solution whenever $\tilde Q$ has a solution. The preceding construction uses exactly 11 unknowns, completing the proof of~\cref{thm:11-unknowns}.
\end{proof}

Using the polynomial $\tilde Q$ it is also possible derive a universal pair, i.e. prove our \cref{thm:11-universal-pair}. We do not formalize this theorem because it depends on a universal pair over $\N$.
\begin{proof}[Proof of \cref{thm:11-universal-pair}]
    Let $(\nu, \delta)_\N$ be a universal pair. Therefore, every polynomial $P$ can be represented in at most this complexity. Theorem~\ref{thm:11-unknowns} establishes that $\tilde Q$ always has eleven integer unknowns, independent of $P$. Calculating the degree of $\tilde Q$ now gives the function $\eta(\nu, \delta)$ as claimed~\cite[Appendix A]{manuscript}. Using the universal pair $(32, 12)_\N$, one can calculate the explicit value for the degree and obtains that 
    \[ (11, 9.19593 \cdot 10^{53})_\Z \]
    is universal. 
\end{proof}

Finally, we present the formal statement of \cref{thm:11-unknowns}. Our formalization implements the combination of \cref{thm:9-unknowns-ZN} and the Three Squares Theorem in an additional 50 lines of Isabelle code. The formal statement assumes \lstinline{is_diophantine_over_N_with A P}, i.e. that $A \subset \N$ is a Diophantine set and represented by the polynomial $P$. Then, by definition, $a \in A$ is equivalent to $\exists z_1, \ldots, z_\nu \in \Z^\nu. \; P(a, z_1, \ldots, z_\nu) = 0$.

\begin{lstlisting}[breaklines=true]
lemma eleven_unknowns_over_Z_polynomial: 
  fixes A :: "nat set" and P :: "int mpoly"
  assumes "is_diophantine_over_N_with A P"
  shows "a $\in$ A = ($\exists$z1 z2 z3 z4 z5 z6 z7 z8 z9 z10 z11.
                  insertion ((!) [int a, z1, z2, z3, z4, z5, z6, z7, z8, z9, z10, z11]) (Q_tilde_poly P) = 0)"
\end{lstlisting}

\section{Conclusion}

We have formalized a universal Diophantine construction that implies a stronger version of Hilbert's Tenth Problem. This construction leads to the novel Theorem~\ref{thm:11-universal-pair} and, moreover, it has been developed largely in parallel to its Isabelle formalization. In the course of the project, we also built a substantial library of advanced facts about multivariate polynomials and formalized results from a textbook on number theory. The total development consists of 20,000 lines of code.

\paragraph*{Experiences from an In-situ Formalization}
The formalization was executed by a student workgroup at a time when only a very early draft of the underlying mathematics was available. Hence we say that this draft manuscript was formalized \emph{in-situ}, in its natural environment---a mathematics department---, and before a version of publication quality had been prepared.

Interacting with a proof assistant significantly shaped the development of the mathematical text: we refined arguments, removed or weakened redundant assumptions, and corrected various minor bugs. For example, we had overlooked a dependency on a lemma in the literature~\cite[Lemma 10]{SunDFI} which only became apparent in the formalization process. We found that for certain technical parts of the proof, Isabelle functioned as a true proof \textit{assistant}: supporting us authors by keeping track of assumptions and bounds so that we could focus on the more conceptual elements of the proof.

\paragraph*{Outlook}
Future research could extend this work by formalizing universal pairs over $\N$. We believe that the infrastructure we have built in this project would provide a strong starting point for such work, eliminating the need to formalize further prerequisites. 

Since this project bridges the frontier between mathematical research and formalization in a proof assistant, further progress will naturally also involve new mathematical discoveries. Hilbert's Tenth Problem has been proven to be solvable for $\delta \leq 2$ or alternatively $\nu = 1$~\cite{9var}, but remains completely open for $\delta = 3$ and largely open for the whole range $2 \leq \nu \leq 10$.

%%
%% Instructions
%%

% \section*{Typesetting instructions -- Summary}
% \label{sec:typesetting-summary}

% Please comply with the following instructions when preparing your article for a LIPIcs proceedings volume. 
% \paragraph*{Minimum requirements}

% \begin{itemize}
% \item Use pdflatex and an up-to-date \LaTeX{} system.
% \item Use further \LaTeX{} packages and custom made macros carefully and only if required.
% \item Use the provided sectioning macros: \verb+\section+, \verb+\subsection+, \verb+\subsubsection+, \linebreak \verb+\paragraph+, \verb+\paragraph*+, and \verb+\subparagraph*+.
% \item Provide suitable graphics of at least 300dpi (preferably in PDF format).
% \item Use BibTeX and keep the standard style (\verb+plainurl+) for the bibliography.
% \item Please try to keep the warnings log as small as possible. Avoid overfull \verb+\hboxes+ and any kind of warnings/errors with the referenced BibTeX entries.
% \item Use a spellchecker to correct typos.
% \end{itemize}

% \paragraph*{Mandatory metadata macros}
% Please set the values of the metadata macros carefully since the information parsed from these macros will be passed to publication servers, catalogues and search engines.
% Avoid placing macros inside the metadata macros. The following metadata macros/environments are mandatory:
% \begin{itemize}
% \item \verb+\title+ and, in case of long titles, \verb+\titlerunning+.
% \item \verb+\author+, one for each author, even if two or more authors have the same affiliation.
% \item \verb+\authorrunning+ and \verb+\Copyright+ (concatenated author names)\\
% The \verb+\author+ macros and the \verb+\Copyright+ macro should contain full author names (especially with regard to the first name), while \verb+\authorrunning+ should contain abbreviated first names.
% \item \verb+\ccsdesc+ (ACM classification, see \url{https://www.acm.org/publications/class-2012}).
% \item \verb+\keywords+ (a comma-separated list of keywords).
% \item \verb+\relatedversion+ (if there is a related version, typically the ``full version''); please make sure to provide a persistent URL, e.\,g., at arXiv.
% \item \verb+\begin{abstract}...\end{abstract}+ .
% \end{itemize}

% \paragraph*{Please do not \ldots} %Do not override the \texttt{\seriesstyle}-defaults}
% Generally speaking, please do not override the \texttt{lipics-v2021}-style defaults. To be more specific, a short checklist also used by Dagstuhl Publishing during the final typesetting is given below.
% In case of \textbf{non-compliance} with these rules Dagstuhl Publishing will remove the corresponding parts of \LaTeX{} code and \textbf{replace it with the \texttt{lipics-v2021} defaults}. In serious cases, we may reject the LaTeX-source and expect the corresponding author to revise the relevant parts.
% \begin{itemize}
% \item Do not use a different main font. (For example, the \texttt{times} package is forbidden.)
% \item Do not alter the spacing of the \texttt{lipics-v2021.cls} style file.
% \item Do not use \verb+enumitem+ and \verb+paralist+. (The \texttt{enumerate} package is preloaded, so you can use
%  \verb+\begin{enumerate}[(a)]+ or the like.)
% \item Do not use ``self-made'' sectioning commands (e.\,g., \verb+\noindent{\bf My+ \verb+Paragraph}+).
% \item Do not hide large text blocks using comments or \verb+\iffalse+ $\ldots$ \verb+\fi+ constructions. 
% \item Do not use conditional structures to include/exclude content. Instead, please provide only the content that should be published -- in one file -- and nothing else.
% \item Do not wrap figures and tables with text. In particular, the package \texttt{wrapfig} is not supported.
% \item Do not change the bibliography style. In particular, do not use author-year citations. (The
% \texttt{natbib} package is not supported.)
% \end{itemize}



%%
%% Bibliography
%%

%% Please use bibtex, 

\bibliography{references}

%\appendix


\end{document}
